%%%%  TEIL II: CHEMIE LESEN UND SPRECHEN  %%%%
\newpage
\part{Chemie lesen und sprechen}

In diesem Teil steht, wie man Elementnamen, Formeln und Reaktionsgleichungen auf Englisch (und zum
Vergleich auf Deutsch) liest, wie Verbindungen benannt werden und wie man Experimente beschreibt.
Links steht der Ausdruck, rechts die englische Lesart und darunter grau die deutsche.

%--------------------------------------------------------------------------------------------------
\section{Die Elemente: Namen und Aussprache}

Viele Elemente heißen im Englischen anders als im Deutschen oder werden deutlich anders
ausgesprochen. Das Elementsymbol ist in beiden Sprachen gleich.

\begin{center}
{\small\renewcommand{\arraystretch}{1.2}
\begin{tabular}{@{}l l l @{\hspace{1cm}} l l l@{}}
	\toprule
	& \textbf{Englisch} & \textbf{Deutsch} & & \textbf{Englisch} & \textbf{Deutsch} \\ \midrule
	H & hydrogen \ipa{"haIdr@dZ@n} & Wasserstoff & K & potassium \ipa{p@"t\ae si@m} & Kalium \\
	He & helium \ipa{"hi:li@m} & Helium & Ca & calcium \ipa{"k\ae lsi@m} & Calcium \\
	Li & lithium \ipa{"lITi@m} & Lithium & Cr & chromium \ipa{"kr@Umi@m} & Chrom \\
	C & carbon \ipa{"kA:b@n} & Kohlenstoff & Mn & manganese \ipa{"m\ae Ng@ni:z} & Mangan \\
	N & nitrogen \ipa{"naItr@dZ@n} & Stickstoff & Fe & iron \ipa{"aI@n} & Eisen \\
	O & oxygen \ipa{"6ksIdZ@n} & Sauerstoff & Cu & copper \ipa{"k6p@} & Kupfer \\
	F & fluorine \ipa{"flO:ri:n} & Fluor & Zn & zinc \ipa{zINk} & Zink \\
	Ne & neon \ipa{"ni:6n} & Neon & Br & bromine \ipa{"br@Umi:n} & Brom \\
	Na & sodium \ipa{"s@Udi@m} & Natrium & Ag & silver \ipa{"sIlv@} & Silber \\
	Mg & magnesium \ipa{m\ae g"ni:zi@m} & Magnesium & Sn & tin \ipa{tIn} & Zinn \\
	Al & aluminium \ipa{""\ae lj@"mIni@m} & Aluminium & I & iodine \ipa{"aI@di:n} & Iod \\
	Si & silicon \ipa{"sIlIk@n} & Silicium & Ba & barium \ipa{"be@ri@m} & Barium \\
	P & phosphorus \ipa{"f6sf@r@s} & Phosphor & W & tungsten \ipa{"t2Nst@n} & Wolfram \\
	S & sulfur \ipa{"s2lf@} & Schwefel & Pt & platinum \ipa{"pl\ae tIn@m} & Platin \\
	Cl & chlorine \ipa{"klO:ri:n} & Chlor & Au & gold \ipa{g@Uld} & Gold \\
	Ar & argon \ipa{"A:g6n} & Argon & Hg & mercury \ipa{"m3:kj@ri} & Quecksilber \\
	 & & & Pb & lead \ipa{led} & Blei \\
	\bottomrule
\end{tabular}}
\end{center}

\anmerkung{Die Halogene enden im Englischen auf \textit{-ine} (\textit{chlorine}), ihre Ionen auf
\textit{-ide} (\textit{chloride}). Die IUPAC empfiehlt die Schreibung \textit{sulfur}; in älteren
britischen Texten steht \textit{sulphur}. Amerikanisch: \textit{aluminum} \ipa{@"lu:mIn@m}. Im Deutschen
gelten heute Iod (früher Jod), Calcium (Kalzium), Oxid (Oxyd) und Ethan (Äthan).}

\subsection*{Beispiele}

\begin{lesen}
\lies{Na}{Sodium is a soft, silvery metal which reacts vigorously with water.}{Natrium ist ein weiches, silbrig glänzendes Metall, das heftig mit Wasser reagiert.}
\lies{Fe}{Iron \ipa{"aI@n} is the most widely used metal. (The r is silent.)}{Eisen ist das am häufigsten verwendete Metall.}
\lies{Pb}{Lead \ipa{led} is a heavy, soft metal -- not to be confused with the verb \textit{to lead} \ipa{li:d}.}{Blei ist ein schweres, weiches Metall.}
\lies{\ce{N2}, \ce{O2}}{Air is a mixture of about seventy-eight per cent nitrogen and twenty-one per cent oxygen.}{Luft ist ein Gemisch aus etwa achtundsiebzig Prozent Stickstoff und einundzwanzig Prozent Sauerstoff.}
\lies{Hg, Br}{Mercury and bromine are the only elements that are liquid at room temperature.}{Quecksilber und Brom sind die einzigen Elemente, die bei Raumtemperatur flüssig sind.}
\end{lesen}

%--------------------------------------------------------------------------------------------------
\section{Formeln und Reaktionsgleichungen lesen}

\begin{center}
\begin{tikzpicture}[font=\small]
	\node[font=\LARGE] (g) at (0,0) {$\color{red}2\,\color{black}\mathrm{H}_{\color{blue}2}\mathrm{O}\,\color{green!50!black}\mathrm{(l)}$};
	\node[align=center] (a) at (-3.6,1.2) {coefficient\\[-2pt]\color{black!60}Koeffizient};
	\node[align=center] (b) at (0,1.6) {chemical symbol\\[-2pt]\color{black!60}Elementsymbol};
	\node[align=center] (c) at (-2.6,-1.3) {subscript\\[-2pt]\color{black!60}Index};
	\node[align=center] (d) at (3.6,1.2) {state symbol\\[-2pt]\color{black!60}Zustandssymbol};
	\draw[-{Stealth}, thin] (a) -- (-1.05,0.2); \draw[-{Stealth}, thin] (b) -- (-0.45,0.35); \draw[-{Stealth}, thin] (c) -- (-0.15,-0.3); \draw[-{Stealth}, thin] (d) -- (1.0,0.25);
\end{tikzpicture}
\end{center}

\begin{lesen}
\lies{\ce{H2O}}{H two O -- water}{H zwei O -- Wasser}
\lies{\ce{CO2}}{C O two -- carbon dioxide}{C O zwei -- Kohlenstoffdioxid}
\lies{\ce{H2SO4}}{H two S O four -- sulfuric acid}{H zwei S O vier -- Schwefelsäure}
\lies{\ce{Ca(OH)2}}{C A, O H twice \;/\; C A, open bracket O H close bracket two -- calcium hydroxide}{C A, O H zweimal \;/\; C A, in Klammern O H, zwei -- Calciumhydroxid}
\lies{\ce{Na+}, \ce{Cl-}}{N A plus -- a sodium ion; C L minus -- a chloride ion}{N A plus -- ein Natrium-Ion; C L minus -- ein Chlorid-Ion}
\lies{\ce{SO4^2-}}{S O four two minus -- the sulfate ion}{S O vier zwei minus -- das Sulfat-Ion}
\lies{(s), (l), (g), (aq)}{state symbols: solid, liquid, gas, aqueous \ipa{"eIkwi@s} (dissolved in water)}{Zustandssymbole: fest, flüssig, gasförmig, in Wasser gelöst}
\lies{\ce{->}}{gives \;/\; yields \;/\; reacts to form}{reagiert zu}
\lies{\ce{<=>}}{is in equilibrium with (a reversible reaction)}{steht im Gleichgewicht mit (umkehrbare Reaktion)}
\lies{\ce{^}, \ce{v}}{a gas is given off; a precipitate is formed}{ein Gas entweicht; ein Niederschlag fällt aus}
\lies{\ce{2 H2 + O2 -> 2 H2O}}{two H two plus O two gives two H two O. \newline In words: Hydrogen reacts with oxygen to form water.}{zwei H zwei plus O zwei reagiert zu zwei H zwei O. \newline In Worten: Wasserstoff reagiert mit Sauerstoff zu Wasser.}
\lies{$\Delta H = -286\eh{kJ/mol}$}{delta H equals minus two hundred and eighty-six kilojoules per mole: the reaction is exothermic}{Delta H gleich minus zweihundertsechsundachtzig Kilojoule pro Mol: Die Reaktion ist exotherm}
\end{lesen}

\subsection*{Beispiele}

\begin{lesen}
\lies{\ce{Zn(s) + 2 HCl(aq) -> ZnCl2(aq) + H2(g)}}{Zinc reacts with hydrochloric acid to form zinc chloride solution and hydrogen gas.}{Zink reagiert mit Salzsäure zu Zinkchloridlösung und gasförmigem Wasserstoff.}
\lies{\ce{CaCO3 -> CaO + CO2 ^}}{On heating, calcium carbonate decomposes into calcium oxide and carbon dioxide.}{Beim Erhitzen zerfällt Calciumcarbonat in Calciumoxid und Kohlenstoffdioxid.}
\lies{\ce{N2 + 3 H2 <=> 2 NH3}}{Nitrogen and hydrogen react to form ammonia; the reaction is reversible.}{Stickstoff und Wasserstoff reagieren zu Ammoniak; die Reaktion ist umkehrbar.}
\lies{\ce{Ba^2+ + SO4^2- -> BaSO4 v}}{Barium ions and sulfate ions form a white precipitate of barium sulfate.}{Barium-Ionen und Sulfat-Ionen bilden einen weißen Niederschlag von Bariumsulfat.}
\lies{Ausgleichen}{Balance the equation: there must be the same number of atoms of each element on both sides.}{Gleiche die Reaktionsgleichung aus: Auf beiden Seiten müssen von jedem Element gleich viele Atome stehen.}
\end{lesen}

%--------------------------------------------------------------------------------------------------
\section{Verbindungen aus zwei Elementen benennen}

Der Name eines Salzes aus zwei Elementen besteht aus dem Namen des Metalls (Kation), bei Bedarf
seiner Wertigkeit als römischer Zahl und dem Namen des Nichtmetalls (Anion) mit der Endung
\textit{-ide} (deutsch: \textit{-id}).

\begin{center}
\begin{tikzpicture}[font=\small]
	\node[font=\LARGE, text=blue] (fe) at (0,0) {Fe};
	\node[font=\LARGE, text=red, anchor=west, inner sep=0pt] (cl) at (0.42,0) {Cl$_3$};
	\node[font=\Large, anchor=west] (n1) at (3.6,0.5) {\textcolor{blue}{iron}\textcolor{green!45!black}{(III)} \textcolor{red}{chloride}};
	\node[font=\Large, anchor=west, text=black!60] (n2) at (3.6,-0.5) {\textcolor{blue!50}{Eisen}\textcolor{green!45!black!60}{(III)}-\textcolor{red!55}{chlorid}};
	\draw[-{Stealth}, thick] (1.6,0) -- (3.3,0);
	\node[align=center, text=blue] (l1) at (2.6,1.9) {metal (cation)\\[-2pt]\color{black!60}Metall (Kation)};
	\node[align=center, text=green!45!black] (l2) at (5.0,1.9) {valency\\[-2pt]\color{black!60}Wertigkeit};
	\node[align=center, text=red] (l3) at (7.6,1.9) {non-metal + \textit{-ide}\\[-2pt]\color{black!60}Nichtmetall + \textit{-id}};
	\draw[thin, blue] (l1.south) -- (4.0,0.8); \draw[thin, green!45!black] (l2.south) -- (4.95,0.8); \draw[thin, red] (l3.south) -- (6.1,0.8);
\end{tikzpicture}
\end{center}

\begin{lesen}
\lies{\ce{NaCl}}{sodium chloride \ipa{"klO:raId}}{Natriumchlorid}
\lies{\ce{MgO}}{magnesium oxide \ipa{"6ksaId}}{Magnesiumoxid}
\lies{\ce{CaF2}}{calcium fluoride \ipa{"flO:raId}}{Calciumfluorid}
\lies{\ce{KBr}, \ce{KI}}{potassium bromide \ipa{"br@UmaId}, potassium iodide \ipa{"aI@daId}}{Kaliumbromid, Kaliumiodid}
\lies{\ce{ZnS}}{zinc sulfide \ipa{"s2lfaId}}{Zinksulfid}
\lies{\ce{FeCl2}, \ce{FeCl3}}{iron(II) chloride, iron(III) chloride -- read: iron two chloride, iron three chloride}{Eisen(II)-chlorid, Eisen(III)-chlorid -- lies: Eisen-zwei-chlorid, Eisen-drei-chlorid}
\lies{\ce{Cu2O}, \ce{CuO}}{copper(I) oxide, copper(II) oxide}{Kupfer(I)-oxid, Kupfer(II)-oxid}
\lies{\ce{PbO2}}{lead(IV) oxide \ipa{""led "fO:r "6ksaId}}{Blei(IV)-oxid}
\end{lesen}

\subsection*{Zahlwörter bei Verbindungen aus Nichtmetallen}

\begin{center}
{\small
\begin{tabular}{@{}l *{10}{c}@{}}
	\toprule
	 & 1 & 2 & 3 & 4 & 5 & 6 & 7 & 8 & 9 & 10 \\ \midrule
	englisch & mono- & di- & tri- & tetra- & penta- & hexa- & hepta- & octa- & nona- & deca- \\
	 & \ipa{"m6n@U} & \ipa{daI} & \ipa{traI} & \ipa{"tetr@} & \ipa{"pent@} & \ipa{"heks@} & \ipa{"hept@} & \ipa{"6kt@} & \ipa{"n6n@} & \ipa{"dek@} \\
	deutsch & mono- & di- & tri- & tetra- & penta- & hexa- & hepta- & octa- & nona- & deca- \\
	\bottomrule
\end{tabular}}
\end{center}

\begin{lesen}
\lies{\ce{CO}}{carbon monoxide \ipa{m@"n6ksaId}}{Kohlenstoffmonoxid}
\lies{\ce{CO2}}{carbon dioxide \ipa{daI"6ksaId}}{Kohlenstoffdioxid}
\lies{\ce{SO2}, \ce{SO3}}{sulfur dioxide, sulfur trioxide \ipa{traI"6ksaId}}{Schwefeldioxid, Schwefeltrioxid}
\lies{\ce{N2O}}{dinitrogen \ipa{daI"naItr@dZ@n} monoxide \;/\; nitrous oxide (laughing gas)}{Distickstoffmonoxid (Lachgas)}
\lies{\ce{NO}, \ce{NO2}}{nitrogen monoxide, nitrogen dioxide}{Stickstoffmonoxid, Stickstoffdioxid}
\lies{\ce{N2O4}}{dinitrogen tetroxide \ipa{te"tr6ksaId}}{Distickstofftetraoxid}
\lies{\ce{N2O5}}{dinitrogen pentoxide \ipa{pen"t6ksaId}}{Distickstoffpentaoxid}
\lies{\ce{HCl}, \ce{H2S}}{hydrogen chloride, hydrogen sulfide}{Chlorwasserstoff, Schwefelwasserstoff}
\lies{\ce{NH3}, \ce{CH4}}{ammonia \ipa{@"m@Uni@}, methane \ipa{"mi:TeIn}}{Ammoniak, Methan}
\end{lesen}

\anmerkung{In älteren englischen Texten stehen statt der römischen Zahlen lateinische Namen: die
Endung \textit{-ous} für die niedrigere, \textit{-ic} für die höhere Wertigkeit -- \textit{ferrous}
\ipa{"fer@s} \textit{chloride} = iron(II) chloride, \textit{ferric} \ipa{"ferIk} \textit{chloride} =
iron(III) chloride, \textit{cuprous} \ipa{"kju:pr@s} und \textit{cupric} \ipa{"kju:prIk} für Kupfer(I) und Kupfer(II).}

\subsection*{Beispiele}

\begin{lesen}
\lies{\ce{Al2O3}}{aluminium oxide -- no Roman numeral is needed, because aluminium always has a valency of three}{Aluminiumoxid -- die römische Zahl entfällt, weil Aluminium immer dreiwertig ist}
\lies{\ce{Fe2O3}}{iron(III) oxide: the main constituent of rust}{Eisen(III)-oxid: der Hauptbestandteil von Rost}
\lies{\ce{P4O10}}{tetraphosphorus decoxide}{Tetraphosphordecaoxid}
\lies{\ce{CCl4}}{carbon tetrachloride \;/\; tetrachloromethane}{Tetrachlormethan}
\lies{\ce{CO}}{Carbon monoxide is a colourless, odourless and very poisonous gas.}{Kohlenstoffmonoxid ist ein farbloses, geruchloses und sehr giftiges Gas.}
\end{lesen}

%--------------------------------------------------------------------------------------------------
\section{Säuren, Basen und ihre Salze benennen}

\begin{center}
\begin{tikzpicture}[font=\small]
	\foreach \y/\a/\s/\da/\ds in {1.6/hydro\ldots\textbf{ic} acid/\ldots\textbf{ide}/\ldots wasserstoffsäure/\ldots\textbf{id},
			0.8/\ldots\textbf{ic} acid/\ldots\textbf{ate}/\ldots säure/\ldots\textbf{at}, 0/\ldots\textbf{ous} acid/\ldots\textbf{ite}/\ldots ige Säure/\ldots\textbf{it}} {
		\node[anchor=east] at (0,\y) {\a}; \draw[-{Stealth}, thick, red] (0.3,\y) -- (1.5,\y); \node[anchor=west] at (1.8,\y) {\s};
		\node[anchor=east, black!60] at (8.6,\y) {\da}; \draw[-{Stealth}, thick, red!50] (8.9,\y) -- (10.1,\y); \node[anchor=west, black!60] at (10.4,\y) {\ds};}
	\node[font=\small\bfseries] at (0.9,2.35) {acid $\rightarrow$ salt}; \node[font=\small\bfseries, black!60] at (9.5,2.35) {Säure $\rightarrow$ Salz};
\end{tikzpicture}
\end{center}

\begin{center}
{\small\renewcommand{\arraystretch}{1.3}
\begin{tabular}{@{}>{\raggedright}p{1.9cm} >{\raggedright}p{4.5cm} >{\raggedright}p{3.7cm} >{\raggedright}p{3.2cm} >{\raggedright\arraybackslash}p{2.1cm}@{}}
	\toprule
	\textbf{Formel} & \textbf{Säure (englisch)} & \textbf{Säure (deutsch)} & \textbf{Salz (englisch)} & \textbf{Salz (deutsch)} \\ \midrule
	\ce{HF} & hydrofluoric acid \ipa{""haIdr@"flO:rIk} & Flusssäure & fluoride & Fluorid \\
	\ce{HCl} & hydrochloric acid \ipa{""haIdr@"kl6rIk} & Salzsäure & chloride & Chlorid \\
	\ce{HBr} & hydrobromic acid \ipa{""haIdr@"br@UmIk} & Bromwasserstoffsäure & bromide & Bromid \\
	\ce{H2S} & hydrogen sulfide (solution) & Schwefelwasserstoff & sulfide & Sulfid \\
	\ce{H2SO4} & sulfuric acid \ipa{s2l"fjU@rIk} & Schwefelsäure & sulfate \ipa{"s2lfeIt} & Sulfat \\
	\ce{H2SO3} & sulfurous acid \ipa{"s2lf@r@s} & schweflige Säure & sulfite \ipa{"s2lfaIt} & Sulfit \\
	\ce{HNO3} & nitric acid \ipa{"naItrIk} & Salpetersäure & nitrate \ipa{"naItreIt} & Nitrat \\
	\ce{HNO2} & nitrous acid \ipa{"naItr@s} & salpetrige Säure & nitrite \ipa{"naItraIt} & Nitrit \\
	\ce{H3PO4} & phosphoric acid \ipa{f6s"f6rIk} & Phosphorsäure & phosphate \ipa{"f6sfeIt} & Phosphat \\
	\ce{H2CO3} & carbonic acid \ipa{kA:"b6nIk} & Kohlensäure & carbonate \ipa{"kA:b@neIt} & Carbonat \\
	\ce{HClO} & hypochlorous acid \ipa{""haIp@"klO:r@s} & hypochlorige Säure & hypochlorite \ipa{""haIp@"klO:raIt} & Hypochlorit \\
	\ce{HClO3} & chloric acid \ipa{"klO:rIk} & Chlorsäure & chlorate \ipa{"klO:reIt} & Chlorat \\
	\ce{HClO4} & perchloric acid \ipa{p@"klO:rIk} & Perchlorsäure & perchlorate \ipa{p@"klO:reIt} & Perchlorat \\
	\ce{CH3COOH} & ethanoic acid \ipa{""eT@"n@UIk} \newline (acetic acid \ipa{@"si:tIk}) & Essigsäure (Ethansäure) & ethanoate \ipa{""eT@"n@UeIt} (acetate \ipa{"\ae sIteIt}) & Acetat (Ethanoat) \\
	\bottomrule
\end{tabular}}
\end{center}

\subsection*{Basen und bekannte Salze}

\begin{lesen}
\lies{\ce{NaOH}}{sodium hydroxide \ipa{haI"dr6ksaId} (caustic soda \ipa{""kO:stIk "s@Ud@}); in water: sodium hydroxide solution}{Natriumhydroxid (Ätznatron); in Wasser gelöst: Natronlauge}
\lies{\ce{KOH}}{potassium hydroxide (caustic potash \ipa{"p6t\ae S}); in water: potassium hydroxide solution}{Kaliumhydroxid (Ätzkali); in Wasser gelöst: Kalilauge}
\lies{\ce{Ca(OH)2}}{calcium hydroxide (slaked lime \ipa{""sleIkt "laIm}); its solution is called limewater}{Calciumhydroxid (gelöschter Kalk); die Lösung heißt Kalkwasser}
\lies{\ce{NH3(aq)}}{ammonia solution (aqueous ammonia)}{Ammoniaklösung (Ammoniakwasser)}
\lies{\ce{NaCl}}{sodium chloride (common salt, table salt)}{Natriumchlorid (Kochsalz)}
\lies{\ce{CaCO3}}{calcium carbonate (limestone, chalk, marble)}{Calciumcarbonat (Kalkstein, Kreide, Marmor)}
\lies{\ce{NaHCO3}}{sodium hydrogencarbonate (sodium bicarbonate, baking soda)}{Natriumhydrogencarbonat (Natron)}
\lies{\ce{KNO3}}{potassium nitrate (saltpetre \ipa{""sO:lt"pi:t@})}{Kaliumnitrat (Kalisalpeter)}
\lies{\ce{CuSO4.5H2O}}{copper(II) sulfate pentahydrate \ipa{""pent@"haIdreIt} (hydrated copper sulfate)}{Kupfer(II)-sulfat-Pentahydrat (blaues Kupfersulfat)}
\end{lesen}

\anmerkung{Für \ce{HCl} gibt es zwei Namen: Das Gas heißt \textit{hydrogen chloride}, seine Lösung in
Wasser \textit{hydrochloric acid}. Die Endung \textit{-ide} wird \ipa{aId} gesprochen, \textit{-ate}
\ipa{eIt} und \textit{-ite} \ipa{aIt} -- \textit{sulfide, sulfate, sulfite} klingen also deutlich verschieden.}

\subsection*{Beispiele}

\begin{lesen}
\lies{\ce{H2SO4 + 2 NaOH -> Na2SO4 + 2 H2O}}{Sulfuric acid reacts with sodium hydroxide to form sodium sulfate and water.}{Schwefelsäure reagiert mit Natriumhydroxid zu Natriumsulfat und Wasser.}
\end{lesen}

%--------------------------------------------------------------------------------------------------
\section{Der pH-Wert und Indikatoren}

\begin{center}
\begin{tikzpicture}[font=\small]
	\phskala{0.9}
	\foreach \k in {0,...,14} {\node[below] at (0.9*\k+0.45,0) {\k};}
	\draw[{Stealth}-, thick] (0,1.6) -- node[above] {acidic -- \color{black!60}sauer} (5.9,1.6);
	\draw[-{Stealth}, thick] (7.6,1.6) -- node[above] {alkaline -- \color{black!60}alkalisch (basisch)} (13.5,1.6);
	\node[above] at (6.75,1.5) {neutral};
	\foreach \k/\e/\d/\y in {1/stomach acid/Magensäure/-0.75, 3/vinegar/Essig/-1.75, 5/rain/Regen/-0.75, 7/pure water/reines Wasser/-1.75, 9/soap solution/Seifenlösung/-0.75, 11/ammonia solution/Ammoniaklösung/-1.75, 13/drain cleaner/Rohrreiniger/-0.75}
		{\draw[thin] (0.9*\k+0.45,-0.45) -- (0.9*\k+0.45,\y); \node[align=center, font=\footnotesize, anchor=north] at (0.9*\k+0.45,\y) {\e\\[-2pt]\color{black!60}\d};}
\end{tikzpicture}
\end{center}

\begin{lesen}
\lies{pH $= 3$}{p H three: the solution is acidic}{p H drei: Die Lösung ist sauer}
\lies{$\mathrm{pH} = -\lg c(\ce{H3O+})$}{pH equals minus log of the concentration of oxonium ions}{pH gleich minus Logarithmus der Konzentration der Oxonium-Ionen}
\lies{Lackmus}{litmus \ipa{"lItm@s}: red in acidic solutions, blue in alkaline solutions}{Lackmus: in saurer Lösung rot, in alkalischer Lösung blau}
\lies{Phenolphthalein}{phenolphthalein \ipa{""fi:n6l"TeIli:n}: colourless in acidic and neutral solutions, pink in alkaline solutions}{Phenolphthalein: in saurer und neutraler Lösung farblos, in alkalischer Lösung rosa bis pink}
\lies{Bromthymolblau}{bromothymol \ipa{""br@Um@U"TaIm6l} blue: yellow in acidic, green in neutral and blue in alkaline solutions}{Bromthymolblau: in saurer Lösung gelb, in neutraler grün und in alkalischer blau}
\end{lesen}

\subsection*{Beispiele}

\begin{lesen}
\lies{$c(\ce{H3O+}) = 0{,}01\eh{mol/l}$ \newline pH $= 2$}{The concentration of oxonium ions is nought point nought one moles per litre, so the pH is two.}{Die Konzentration der Oxonium-Ionen beträgt null Komma null eins Mol pro Liter, der pH-Wert ist also zwei.}
\lies{pH $3 \rightarrow$ pH $2$}{A decrease of one pH unit means that the solution is ten times more acidic.}{Sinkt der pH-Wert um eine Einheit, ist die Lösung zehnmal so sauer.}
\end{lesen}

%--------------------------------------------------------------------------------------------------
\section{Organische Verbindungen benennen}

Der Stamm des Namens gibt die Zahl der Kohlenstoffatome an, die Endung die Stoffklasse.

\begin{center}
{\small\renewcommand{\arraystretch}{1.15}
\begin{tabular}{@{}c l l l @{\hspace{1cm}} c l l l@{}}
	\toprule
	\textbf{C} & \textbf{Formel} & \textbf{Englisch} & \textbf{Deutsch} & \textbf{C} & \textbf{Formel} & \textbf{Englisch} & \textbf{Deutsch} \\ \midrule
	1 & \ce{CH4} & methane \ipa{"mi:TeIn} & Methan & 6 & \ce{C6H14} & hexane \ipa{"hekseIn} & Hexan \\
	2 & \ce{C2H6} & ethane \ipa{"i:TeIn} & Ethan & 7 & \ce{C7H16} & heptane \ipa{"hepteIn} & Heptan \\
	3 & \ce{C3H8} & propane \ipa{"pr@UpeIn} & Propan & 8 & \ce{C8H18} & octane \ipa{"6kteIn} & Octan \\
	4 & \ce{C4H10} & butane \ipa{"bju:teIn} & Butan & 9 & \ce{C9H20} & nonane \ipa{"n@UneIn} & Nonan \\
	5 & \ce{C5H12} & pentane \ipa{"penteIn} & Pentan & 10 & \ce{C10H22} & decane \ipa{"dekeIn} & Decan \\
	\bottomrule
\end{tabular}}
\end{center}

\begin{center}
{\footnotesize\renewcommand{\arraystretch}{1.25}
\begin{tabular}{@{}l l l l l@{}}
	\toprule
	\textbf{Stoffklasse (englisch)} & \textbf{deutsch} & \textbf{Endung} & \textbf{Beispiel (englisch)} & \textbf{deutsch} \\ \midrule
	alkanes & Alkane & -ane \,/\, -an & ethane & Ethan \\
	alkenes & Alkene & -ene \,/\, -en & ethene \ipa{"eTi:n} & Ethen \\
	alkynes & Alkine & -yne \,/\, -in & ethyne \ipa{"i:TaIn} & Ethin \\
	alcohols & Alkohole (Alkanole) & -ol \,/\, -ol & ethanol \ipa{"eT@n6l} & Ethanol \\
	aldehydes & Aldehyde (Alkanale) & -al \,/\, -al & ethanal \ipa{"eT@n\ae l} & Ethanal \\
	ketones & Ketone (Alkanone) & -one \,/\, -on & propanone & Propanon \\
	carboxylic acids & Carbonsäuren & -oic acid \,/\, -säure & ethanoic acid & Ethansäure \\
	esters & Ester & -yl -oate \,/\, -ester & ethyl ethanoate & Ethansäureethylester \\
	\bottomrule
\end{tabular}}
\end{center}

\begin{lesen}
\lies{\struktur{C(-[:120]H)(-[:240]H)=C(-[:60]H)(-[:-60]H)}}{ethene (common name: ethylene \ipa{"eT@li:n})}{Ethen (Trivialname: Ethylen)}
\lies{\struktur{H-C~C-H}}{ethyne (common name: acetylene \ipa{@"set@li:n})}{Ethin (Trivialname: Acetylen)}
\lies{\struktur{H_3C-CH(-[6]CH_3)-CH_3}}{2-methylpropane -- read: two methyl \ipa{"meTIl} propane (common name: isobutane)}{2-Methylpropan -- lies: zwei Methylpropan (Trivialname: Isobutan)}
\lies{\struktur{H_3C-CH_2-OH}}{ethanol (ethyl alcohol)}{Ethanol (Ethylalkohol, Trinkalkohol)}
\lies{\struktur{HO-CH_2-CH(-[6]OH)-CH_2-OH}}{propane-1,2,3-triol -- read: propane one two three triol (common name: glycerol \ipa{"glIs@r6l})}{Propan-1,2,3-triol (Trivialname: Glycerin)}
\lies{\struktur{H-C(=[2]O)-H}}{methanal (common name: formaldehyde \ipa{fO:"m\ae ldIhaId})}{Methanal (Trivialname: Formaldehyd)}
\lies{\struktur{H_3C-C(=[2]O)-CH_3}}{propanone (common name: acetone \ipa{"\ae sIt@Un})}{Propanon (Trivialname: Aceton)}
\lies{\struktur{H-C(=[2]O)-OH}}{methanoic acid \ipa{""meT@"n@UIk} (common name: formic \ipa{"fO:mIk} acid)}{Methansäure (Trivialname: Ameisensäure)}
\lies{\struktur{H_3C-C(=[2]O)-OH}}{ethanoic acid (common name: acetic acid)}{Ethansäure (Trivialname: Essigsäure)}
\lies{\struktur{H_3C-C(=[2]O)-O-CH_2-CH_3}}{ethyl ethanoate (common name: ethyl acetate)}{Ethansäureethylester (Trivialname: Essigsäureethylester)}
\lies{\struktur{*6(-=-=-=)}}{benzene \ipa{"benzi:n}}{Benzol (Benzen)}
\lies{\struktur{*6(-=-=(-CH_3)-=)}}{methylbenzene (common name: toluene \ipa{"t6lju:i:n})}{Methylbenzol (Trivialname: Toluol)}
\lies{\struktur{*6(-=-=(-OH)-=)}}{phenol \ipa{"fi:n6l}}{Phenol (Hydroxybenzol)}
\end{lesen}

\anmerkung{Falsche Freunde: Benzol heißt auf Englisch \textit{benzene}; das deutsche Benzin ist
\textit{petrol} (britisch) bzw.\ \textit{gasoline} (amerikanisch). Das englische \textit{benzine}
bezeichnet Waschbenzin. Stellungsziffern liest man einzeln: \textit{but-1-ene} -- „but one ene“,
\textit{propan-2-ol} -- „propan two ol“, \textit{2,2-dimethylpropane} -- „two two dimethylpropane“.}

\subsection*{Beispiele}

\begin{lesen}
\lies{\ce{C3H8}}{Propane is the third member of the homologous series of the alkanes.}{Propan ist das dritte Glied der homologen Reihe der Alkane.}
\lies{\ce{CH3-CH=CH2}}{propene: an alkene with three carbon atoms}{Propen: ein Alken mit drei Kohlenstoffatomen}
\lies{\ce{CH3-CH(OH)-CH3}}{propan-2-ol: the hydroxyl group is attached to the second carbon atom}{Propan-2-ol: Die Hydroxygruppe sitzt am zweiten Kohlenstoffatom}
\lies{\ce{CH2Cl-CH2Cl}}{1,2-dichloroethane -- read: one two dichloro ethane}{1,2-Dichlorethan -- lies: eins zwei Dichlorethan}
\lies{\ce{C2H5OH + CH3COOH <=> CH3COOC2H5 + H2O}}{Ethanol reacts with ethanoic acid to form ethyl ethanoate and water (esterification).}{Ethanol reagiert mit Ethansäure zu Ethansäureethylester und Wasser (Veresterung).}
\end{lesen}

%--------------------------------------------------------------------------------------------------
\section{Rechnen mit Stoffmengen}

\begin{lesen}
\lies{$n = \dfrac{m}{M}$}{$n$ equals $m$ over capital $M$: the amount of substance equals the mass divided by the molar mass}{$n$ gleich $m$ durch groß $M$: Die Stoffmenge ist die Masse durch die molare Masse}
\lies{$N = n \cdot N_A$}{the number of particles equals the amount of substance times the Avogadro constant}{die Teilchenanzahl ist die Stoffmenge mal die Avogadro-Konstante}
\lies{$c = \dfrac{n}{V}$}{$c$ equals $n$ over $V$: the concentration equals the amount of substance divided by the volume of the solution}{$c$ gleich $n$ durch $V$: Die Konzentration ist die Stoffmenge durch das Volumen der Lösung}
\lies{$V = n \cdot V_m$}{the volume of a gas equals the amount of substance times the molar volume}{das Volumen eines Gases ist die Stoffmenge mal das molare Volumen}
\lies{$c(\ce{HCl}) = 0{,}1\eh{mol/l}$}{the concentration of the hydrochloric acid is nought point one moles per litre (nought point one molar)}{die Konzentration der Salzsäure beträgt null Komma eins Mol pro Liter}
\lies{$M(\ce{NaCl}) = 58{,}5\eh{g/mol}$}{the molar mass of sodium chloride is fifty-eight point five grams per mole}{die molare Masse von Natriumchlorid beträgt achtundfünfzig Komma fünf Gramm pro Mol}
\end{lesen}

\anmerkung{Im Englischen schreibt man die Konzentration meist in $\mathrm{mol/dm^3}$ oder kurz M
(\textit{molar}): $0{,}1\eh{mol/l} = 0.1\eh{mol\,dm^{-3}} = 0.1\eh{M}$. Eckige Klammern stehen für
Konzentrationen: $[\ce{H3O+}]$.}

\subsection*{Beispiele}

\begin{lesen}
\lies{$M(\ce{CO2}) = (12 + 2 \cdot 16)\eh{g/mol}$ \newline $M(\ce{CO2}) = 44\eh{g/mol}$}{To find the molar mass, add up the atomic masses: twelve plus two times sixteen is forty-four grams per mole.}{Für die molare Masse addiert man die Atommassen: zwölf plus zwei mal sechzehn ergibt vierundvierzig Gramm pro Mol.}
\lies{\ce{2 Mg + O2 -> 2 MgO} \newline $n(\ce{Mg}) = 1\eh{mol}$}{What mass of magnesium oxide is formed when twenty-four point three grams of magnesium burn? First calculate the amount of magnesium: one mole.}{Welche Masse Magnesiumoxid entsteht, wenn vierundzwanzig Komma drei Gramm Magnesium verbrennen? Berechne zuerst die Stoffmenge des Magnesiums: ein Mol.}
\lies{$n(\ce{MgO}) = n(\ce{Mg}) = 1\eh{mol}$ \newline $m(\ce{MgO}) = 1\eh{mol} \cdot 40{,}3\eh{g/mol}$ \newline $m(\ce{MgO}) = 40{,}3\eh{g}$}{From the equation, two moles of magnesium give two moles of magnesium oxide. So forty point three grams of magnesium oxide are formed.}{Nach der Reaktionsgleichung entstehen aus zwei Mol Magnesium zwei Mol Magnesiumoxid. Es entstehen also vierzig Komma drei Gramm Magnesiumoxid.}
\lies{$m = c \cdot V \cdot M$ \newline $m = 0{,}1\eh{\tfrac{mol}{l}} \cdot 1\eh{l} \cdot 40\eh{\tfrac{g}{mol}} = 4\eh{g}$}{To make one litre of sodium hydroxide solution with a concentration of nought point one moles per litre, dissolve four grams of sodium hydroxide in water and make the solution up to one litre.}{Für einen Liter Natronlauge der Konzentration null Komma eins Mol pro Liter löst man vier Gramm Natriumhydroxid in Wasser und füllt auf einen Liter auf.}
\lies{$c(\ce{HCl}) = \dfrac{0{,}1\eh{mol/l} \cdot 25\eh{ml}}{20\eh{ml}}$ \newline $c(\ce{HCl}) = 0{,}125\eh{mol/l}$}{Twenty millilitres of hydrochloric acid are neutralised by twenty-five millilitres of sodium hydroxide solution of concentration nought point one moles per litre. The concentration of the acid is nought point one two five moles per litre.}{Zwanzig Milliliter Salzsäure werden von fünfundzwanzig Millilitern Natronlauge der Konzentration null Komma eins Mol pro Liter neutralisiert. Die Konzentration der Säure beträgt null Komma eins zwei fünf Mol pro Liter.}
\end{lesen}

%--------------------------------------------------------------------------------------------------
\section{Experimente beschreiben und sicher arbeiten}

\begin{center}
\begin{tikzpicture}[font=\small, line cap=round, line join=round]
	\begin{scope}[scale=1.15]
		\brenner{0}{0}
		\draw[black, very thick] (-1.3,2.6) -- (1.3,2.6) (-1.15,2.6) -- (-1.5,0) (1.15,2.6) -- (1.5,0);
		\draw[black!60, thick, densely dashed] (-1.4,2.72) -- (1.4,2.72);
		\fill[blue!12] (-0.9,2.78) rectangle (0.9,3.8); \draw[blue, very thick] (-1.05,4.7) -- (-0.9,4.6) -- (-0.9,2.78) -- (0.9,2.78) -- (0.9,4.7);
		\draw[blue, thick, rounded corners=1.5pt] (0.3,3.1) rectangle (0.48,5.3); \fill[red] (0.39,3.1) circle (0.13); \fill[red] (0.36,3.1) rectangle (0.42,4.2);
	\end{scope}
	\foreach \y/\e/\d/\px/\py in {5.4/thermometer/Thermometer/0.6/5.6, 4.5/beaker/Becherglas/1.1/4.6, 3.6/gauze/Drahtnetz/1.5/3.15, 2.7/tripod/Dreifuß/1.6/1.9, 1.8/Bunsen burner/Gasbrenner/0.3/1.0, 0.9/gas supply/Gaszufuhr/-1.1/0.15}
		{\node[anchor=west] (n) at (4,\y) {\e{} -- \color{black!60}\d}; \draw[-{Stealth}, thin] (n.west) -- (\px,\py);}
\end{tikzpicture}
\end{center}

Die Teile eines Versuchsprotokolls (\textit{lab report}):

\begin{lesen}
\lies{\textbf{Aufgabe}}{\textbf{Aim:} To investigate how temperature affects the solubility of potassium nitrate in water.}{Untersuche, wie die Temperatur die Löslichkeit von Kaliumnitrat in Wasser beeinflusst.}
\lies{\textbf{Geräte und Chemikalien}}{\textbf{Apparatus and chemicals:} beaker, tripod, gauze, Bunsen burner, thermometer, balance; potassium nitrate, distilled water}{Becherglas, Dreifuß, Drahtnetz, Gasbrenner, Thermometer, Waage; Kaliumnitrat, destilliertes Wasser}
\lies{\textbf{Durchführung}}{\textbf{Method:} Pour fifty millilitres of water into the beaker. Add potassium nitrate a spatula at a time and stir until no more dissolves. Heat the mixture and repeat.}{Gib fünfzig Milliliter Wasser in das Becherglas. Füge spatelweise Kaliumnitrat hinzu und rühre, bis sich nichts mehr löst. Erwärme das Gemisch und wiederhole den Vorgang.}
\lies{\textbf{Beobachtung}}{\textbf{Observations:} More potassium nitrate dissolves in hot water than in cold water. On cooling, crystals form.}{In heißem Wasser löst sich mehr Kaliumnitrat als in kaltem. Beim Abkühlen bilden sich Kristalle.}
\lies{\textbf{Auswertung}}{\textbf{Conclusion:} The solubility of potassium nitrate increases with temperature.}{Die Löslichkeit von Kaliumnitrat nimmt mit steigender Temperatur zu.}
\end{lesen}

\subsection*{Beobachtungen und Nachweise}

\begin{lesen}
\lies{Gasentwicklung}{Bubbles of gas are given off (effervescence \ipa{""ef@"ves=ns}).}{Es steigen Gasblasen auf (die Lösung schäumt).}
\lies{Farbänderung}{The solution turns from blue to colourless.}{Die Lösung entfärbt sich von Blau nach Farblos.}
\lies{Niederschlag}{A white precipitate forms.}{Es bildet sich ein weißer Niederschlag.}
\lies{Sauerstoff}{Test for oxygen: a glowing splint relights.}{Nachweis von Sauerstoff (Glimmspanprobe): Ein glimmender Holzspan flammt auf.}
\lies{Wasserstoff}{Test for hydrogen: a lighted splint gives a squeaky pop.}{Nachweis von Wasserstoff (Knallgasprobe): Ein brennender Span erzeugt ein pfeifendes Geräusch.}
\lies{Kohlenstoffdioxid}{Test for carbon dioxide: limewater turns milky.}{Nachweis von Kohlenstoffdioxid: Kalkwasser trübt sich.}
\lies{Flammenfärbung}{Flame test: sodium compounds colour the flame yellow, potassium compounds lilac.}{Flammenfärbung: Natriumverbindungen färben die Flamme gelb, Kaliumverbindungen violett.}
\end{lesen}

\subsection*{Gefahren und Sicherheit}

\begin{lesen}
\lies{entzündbar}{flammable \ipa{"fl\ae m@b=l}}{}
\lies{ätzend}{corrosive \ipa{k@"r@UsIv}}{}
\lies{giftig}{toxic \ipa{"t6ksIk}}{}
\lies{gesundheitsschädlich, reizend}{harmful, irritant \ipa{"IrIt@nt}}{}
\lies{brandfördernd}{oxidising \ipa{"6ksIdaIzIN}}{}
\lies{explosiv}{explosive \ipa{Ik"spl@UsIv}}{}
\lies{umweltgefährlich}{hazardous to the environment \ipa{"h\ae z@d@s}}{}
\end{lesen}

\subsection*{Beispiele}

\begin{lesen}
\lies{Schutz}{Wear safety goggles and tie back long hair.}{Trage eine Schutzbrille und binde lange Haare zusammen.}
\lies{Verdünnen}{Always add acid to water, never water to acid.}{Gib immer die Säure ins Wasser, nie das Wasser in die Säure.}
\lies{Geruchsprobe}{To smell a substance, waft the vapour towards your nose with your hand.}{Fächle dir zur Geruchsprobe die Dämpfe mit der Hand zu.}
\lies{Abzug}{Carry out the experiment in a fume cupboard.}{Führe den Versuch unter dem Abzug durch.}
\lies{Entsorgung}{Do not pour the residues down the sink; put them in the waste container.}{Gieße die Reste nicht in den Ausguss; gib sie in den Sammelbehälter.}
\end{lesen}
